\documentclass[letterpaper, italian]{report}
\usepackage[utf8]{inputenc}
\usepackage{csvsimple}
\usepackage[hidelinks, unicode]{hyperref}
\usepackage{graphicx}
\usepackage[T1]{fontenc}
\usepackage{contour}
\usepackage{ulem}
\usepackage{bookmark}
\usepackage{enumitem}
\usepackage{longtable}
\usepackage{booktabs}
\usepackage{float}
\usepackage{xcolor}
% \usepackage[landscape,twocolumn]{geometry} here if we want to make it landscape later

\renewcommand{\ULdepth}{1.8pt}
\contourlength{0.8pt}

\newcommand{\myuline}[1]{%
  \uline{\phantom{#1}}%
  \llap{\contour{white}{#1}}%
}

\renewcommand{\thesection}{\arabic{section}} % removes reference of \chapter to avoid "0."
\renewcommand{\thesubsection}{\thesection.\alph{subsection}} % removes reference of \chapter to avoid "0."
% \titleformat{\chapter}{\normalfont\huge}{\thechapter}{20pt}{\huge\it}

\title{
    \leavevmode{\includegraphics[width=0.8\textwidth]{resources/Universita-degli-studi-di-torino-logo.png}\newline\newline}\\
    Cat \& Ring \\
    \large Gestire le Ricette (e le Preparazioni)
}
\author{Eduard Antonovic Occhipinti, Riccardo Cardona}

\begin{document}

\maketitle

\tableofcontents

\part{Requisiti}
\chapter{UC Dettagliato}

\section*{Informazioni generali}
\textbf{Nome caso d'uso}{: Gestire le ricette}\newline
\textbf{Portata}{: Sistema}\newline
\textbf{Livello}{: Obiettivo utente}\newline
\textbf{Attore primario}{: Cuoco}\newline
\textbf{Parti Interessate}\newline
\textbf{Pre-condizioni}{: L'attore deve essere autenticato come Cuoco}\newline
\textbf{Garanzie di successo o post-condizioni}{: La procedura di cucina è salvata}

\section*{Scenario Principale di Successo}

\begin{table}[H]\centering
\csvreader[
    tabular = |c|p{7cm}|p{4cm}|,
    respect sharp,
    table head = \hline\bfseries \# & \bfseries Attore & \bfseries Sistema\\\hline,
    late after line = \\\hline
]{resources/csv/scenario_principale_di_successo.csv}{}{
    \csvcoli & \csvcolii & \csvcoliii
}
\end{table}

\section*{Estensioni 1}
\begin{table}[H]\centering\caption{Estensione 1a}
\csvreader[
    tabular = |c|p{7cm}|p{4cm}|,
    respect sharp,
    table head = \hline\bfseries \# & \bfseries Attore & \bfseries Sistema\\\hline,
    late after line = \\\hline
]{resources/csv/estensione1a.csv}{}{
    \csvcoli & \csvcolii & \csvcoliii
}
\end{table}

\begin{table}[H]\centering\caption{Estensione 1b}
\csvreader[
      tabular = |c|p{7cm}|p{4cm}|,
      respect sharp,
      table head = \hline\bfseries \# & \bfseries Attore & \bfseries Sistema\\\hline,
      late after line = \\\hline
]{resources/csv/estensione1b.csv}{}{
      \csvcoli & \csvcolii & \csvcoliii
}
\end{table}

\begin{table}[H]\centering\caption{Estensione 1c}
\csvreader[
      tabular = |c|p{7cm}|p{4cm}|,
      respect sharp,
      table head = \hline\bfseries \# & \bfseries Attore & \bfseries Sistema\\\hline,
      late after line = \\\hline
]{resources/csv/estensione1c.csv}{}{
      \csvcoli & \csvcolii & \csvcoliii
}
\end{table}

\section*{Eccezioni 1}
\begin{table}[H]\centering\color{red}\caption{Eccezione 1b.1a}
\csvreader[
      tabular = |c|p{7cm}|p{4cm}|,
      respect sharp,
      table head = \hline\bfseries \# & \bfseries Attore & \bfseries Sistema\\\hline,
      late after line = \\\hline
]{resources/csv/eccezione1b1a.csv}{}{
      \csvcoli & \csvcolii & \csvcoliii
}
\end{table}

\begin{table}[H]\centering\color{red}\caption{Eccezione 1b.1b}
\csvreader[
      tabular = |c|p{7cm}|p{4cm}|,
      respect sharp,
      table head = \hline\bfseries \# & \bfseries Attore & \bfseries Sistema\\\hline,
      late after line = \\\hline
]{resources/csv/eccezione1b1b.csv}{}{
      \csvcoli & \csvcolii & \csvcoliii
}
\end{table}

\begin{table}[H]\centering\color{red}\caption{Eccezione 1c.1a}
\csvreader[
      tabular = |c|p{7cm}|p{4cm}|,
      respect sharp,
      table head = \hline\bfseries \# & \bfseries Attore & \bfseries Sistema\\\hline,
      late after line = \\\hline
]{resources/csv/eccezione1c1a.csv}{}{
      \csvcoli & \csvcolii & \csvcoliii
}
\end{table}

\begin{table}[H]\centering\color{red}\caption{Eccezione 1c.1b}
\csvreader[
      tabular = |c|p{7cm}|p{4cm}|,
      respect sharp,
      table head = \hline\bfseries \# & \bfseries Attore & \bfseries Sistema\\\hline,
      late after line = \\\hline
]{resources/csv/eccezione1c1b.csv}{}{
      \csvcoli & \csvcolii & \csvcoliii
}
\end{table}

\section*{Estensioni 2}
\begin{table}[H]\centering\caption{Estensione 2a}
\csvreader[
      tabular = |c|p{7cm}|p{4cm}|,
      respect sharp,
      table head = \hline\bfseries \# & \bfseries Attore & \bfseries Sistema\\\hline,
      late after line = \\\hline
]{resources/csv/estensione2a.csv}{}{
      \csvcoli & \csvcolii & \csvcoliii
}
\end{table}

\section*{Estensioni 3}
\begin{table}[H]\centering\caption{Estensione 3a}
\csvreader[
      tabular = |c|p{7cm}|p{4cm}|,
      respect sharp,
      table head = \hline\bfseries \# & \bfseries Attore & \bfseries Sistema\\\hline,
      late after line = \\\hline
]{resources/csv/estensione3a.csv}{}{
      \csvcoli & \csvcolii & \csvcoliii
}
\end{table}

\begin{table}[H]\centering\caption{Estensione 3b}
\csvreader[
      tabular = |c|p{7cm}|p{4cm}|,
      respect sharp,
      table head = \hline\bfseries \# & \bfseries Attore & \bfseries Sistema\\\hline,
      late after line = \\\hline
]{resources/csv/estensione3b.csv}{}{
      \csvcoli & \csvcolii & \csvcoliii
}
\end{table}

\section*{Estensioni 5}
\begin{table}[H]\centering\caption{Estensione 5a}
\csvreader[
      tabular = |c|p{7cm}|p{4cm}|,
      respect sharp,
      table head = \hline\bfseries \# & \bfseries Attore & \bfseries Sistema\\\hline,
      late after line = \\\hline
]{resources/csv/estensione5a.csv}{}{
      \csvcoli & \csvcolii & \csvcoliii
}
\end{table}

\section*{Estensioni 6}
\begin{table}[H]\centering\caption{Estensione 6a}
\csvreader[
      tabular = |c|p{7cm}|p{4cm}|,
      respect sharp,
      table head = \hline\bfseries \# & \bfseries Attore & \bfseries Sistema\\\hline,
      late after line = \\\hline
]{resources/csv/estensione6a.csv}{}{
      \csvcoli & \csvcolii & \csvcoliii
}
\end{table}

\begin{table}[H]\centering\caption{Estensione 6b}
\csvreader[
      tabular = |c|p{7cm}|p{4cm}|,
      respect sharp,
      table head = \hline\bfseries \# & \bfseries Attore & \bfseries Sistema\\\hline,
      late after line = \\\hline
]{resources/csv/estensione6b.csv}{}{
      \csvcoli & \csvcolii & \csvcoliii
}
\end{table}

\section*{Estensioni 7}
\begin{table}[H]\centering\caption{Estensione 7a}
\csvreader[
      tabular = |c|p{7cm}|p{4cm}|,
      respect sharp,
      table head = \hline\bfseries \# & \bfseries Attore & \bfseries Sistema\\\hline,
      late after line = \\\hline
]{resources/csv/estensione7a.csv}{}{
      \csvcoli & \csvcolii & \csvcoliii
}
\end{table}

\begin{table}[H]\centering\caption{Estensione 7b}
\csvreader[
      tabular = |c|p{7cm}|p{4cm}|,
      respect sharp,
      table head = \hline\bfseries \# & \bfseries Attore & \bfseries Sistema\\\hline,
      late after line = \\\hline
]{resources/csv/estensione7b.csv}{}{
      \csvcoli & \csvcolii & \csvcoliii
}
\end{table}

\begin{table}[H]\centering\caption{Estensione 7c}
\csvreader[
      tabular = |c|p{7cm}|p{4cm}|,
      respect sharp,
      table head = \hline\bfseries \# & \bfseries Attore & \bfseries Sistema\\\hline,
      late after line = \\\hline
]{resources/csv/estensione7c.csv}{}{
      \csvcoli & \csvcolii & \csvcoliii
}
\end{table}

\section*{Estensioni 8}
\begin{table}[H]\centering\caption{Estensione 8a}
\csvreader[
      tabular = |c|p{7cm}|p{4cm}|,
      respect sharp,
      table head = \hline\bfseries \# & \bfseries Attore & \bfseries Sistema\\\hline,
      late after line = \\\hline
]{resources/csv/estensione8a.csv}{}{
      \csvcoli & \csvcolii & \csvcoliii
}
\end{table}

\chapter{Modello di Dominio e SSD}
\section{Modello di Dominio}
\section{Diagrammi di Sequenza di Sistema}

\chapter{Contratti delle Operazioni}
\textbf{Pre-condizione generale}{: l'utente è identificato come Cuoco o come Chef}\textit{c}

\section[creaProcedura]{\texorpdfstring{creaProcedura(\myuline{nome}: Testo, \myuline{èRicetta}: sì/no)}}\label{h.n5k8dzisnnko}
\textbf{Pre-condizioni:} -\newline
\textbf{Post-condizioni:}
\begin{itemize}[label=$-$, noitemsep]
      \item {[\textbf{Se} \myuline{èRicetta} == sì]}
      \setlist{nolistsep}
      \begin{itemize}[label=$-$, noitemsep]
            \item È stata creata un'istanza \textit{p} di Ricetta
            \item \textit{p.pubblicato} = no
            \item \textit{p.nome } = \myuline{nome}
      \end{itemize}
      \item {[\textbf{Se} \myuline{èRicetta} == no]}
      \setlist{nolistsep}
      \begin{itemize}[label=$-$, noitemsep]
            \item È stata creata un'istanza \textit{p} di Preparazione
            \item \textit{p.pubblicato} = no
            \item \textit{p.nome } = \myuline{nome}
      \end{itemize}
      \item \textit{p}\textbf{ si trova} nel ricettario
      \item \textit{c }\textbf{è proprietario di}\textit{ p}
\end{itemize}

\newpage
\subsection[creaVarianteProcedura]{\texorpdfstring{creaVarianteProcedura(\myuline{procedura}: Procedura, \myuline{}{nome}: Testo)}}\label{h.4pfv8u235m4y}
\textbf{Pre-condizioni:} -\newline
\textbf{Post-condizioni:}
\begin{itemize}[label=$-$, noitemsep]
      \item {È stata creata un'istanza }\textit{p}{ di Procedura}
      \setlist{nolistsep}
      \begin{itemize}[label=$-$, noitemsep]
            \item \textit{p.pubblicato }{= no}
            \item \textit{p.tempo }{= }\myuline{procedura.tempo}
            \item \textit{p.nome }{= }\myuline{}{nome}
            \item \textit{p.preparabileInAnticipo }{= }\myuline{procedura.preparabileInAnticipo}
            \item {Per ogni Passo }\textit{pas}{ tale che }\myuline{procedura}
                  \textbf{ è formata da }\textit{pas }{è stata creata un'istanza }
                  \textit{ps}{ di Passo tale che}
            \setlist{nolistsep}
            \begin{itemize}[label=$-$, noitemsep]
                  \item \textit{ps.descrizione }{= }\textit{pas.descrizione}
                  \item \textit{p }{è}{~}\textbf{formata da }\textit{ps}
            \end{itemize}
            \item {Per ogni Ingrediente }\textit{ing}{~tale che }\textit{ing}{~}
                  {compone}{~}\myuline{procedura}{~}{è stata creata un'istanza }
                  \textit{i}{~di Ingrediente tale che}
            \setlist{nolistsep}
            \begin{itemize}[label=$-$, noitemsep]
                  \item \textit{i.nome }{= }\textit{ing.nome}
                  \item \textit{i.dose }{= }\textit{ing.dose}
                  \item \textit{i }\textbf{compone}\textit{~p}
            \end{itemize}
      \item \textit{p }\textbf{variante di }\myuline{procedura}
      \item \textit{p }\textbf{si trova nel }{Ricettario}
      \end{itemize}
      \item \textit{c }\textbf{è proprietario di}\textit{~p}
\end{itemize}

\subsection[apriProceduraPerModifica]{\texorpdfstring{apriProceduraPerModifica(\myuline{procedura}: Procedura)}}\label{h.y0wiqdm1uq8i}
\textbf{Pre-condizioni: }{-}\newline
\textbf{Post-condizioni: }
\begin{itemize}[label=$-$, noitemsep]
      \item {se c }\textbf{è proprietario di}{~}\myuline{procedura}{, }\myuline{procedura}{~non
        }\textbf{riguarda}{~nessun Assegnamento e non }\textbf{serve }{per nessun Menù}
      \setlist{nolistsep}
      \begin{itemize}[label=$-$, noitemsep]
            \item {{[}}\textbf{Se}{~}\myuline{procedura.pubblicato}{~= si{]} }\myuline{procedura.pubblicato}{~= no}
      \end{itemize}
\end{itemize}


\subsection[eliminaProcedura]{\texorpdfstring{eliminaProcedura(\myuline{procedura}: Procedura)}}\label{h.au1glj4lcszr}
\textbf{Pre-condizioni: }{-}\newline
\textbf{Post-condizioni:}
\begin{itemize}[label=$-$, noitemsep]
      \item {se }\textit{c }{è }\textbf{proprietario di}{~}\myuline{procedura}{, }\myuline{procedura}{~non }
      \textbf{riguarda}{~nessun Assegnamento e non }\textbf{serve }{per nessun Menù}
      \setlist{nolistsep}
      \begin{itemize}[label=$-$, noitemsep]
            \item {per ogni Passo }\textit{p}{ tale che }\myuline{procedura}{~è }
                  \textbf{formata}{~da }\textit{p}{, }\textit{p}{ è eliminato}
            \item {per ogni Ingrediente }\textit{i}{ tale che }\myuline{procedura}
                  {~è }\textbf{composto}{~da }\textit{i}{, }\textit{i}{ è eliminato}
      \end{itemize}
      \item \myuline{procedura}{~è eliminata}
\end{itemize}

\newpage
\section[aggiungePasso]{\texorpdfstring{aggiungePasso(\myuline{descrizione}:{~Testo, }\myuline{prep}: sì/no)}}\label{h.tszsgr5cfwm3}
\textbf{Pre-condizioni:}
\begin{itemize}[label=$-$, noitemsep]
      \item {è in corso la definizione di una Procedura }\textit{p}
\end{itemize}
\textbf{Post-condizioni:}
\begin{itemize}[label=$-$, noitemsep]
      \item {è stata creata una istanza }\textit{pa}{ di Passo}
      \item \textit{pa.descrizione = }\myuline{descrizione}
      \item \textit{pa.preparabileInAnticipo }{= }\myuline{prep}
      \item \textit{p}{ è }\textbf{formato }{da }\textit{pa}
\end{itemize}

\subsection[creaRaggruppamento]{\texorpdfstring{creaRaggruppamento (\myuline{descrizione}: testo)}}\label{a.1-crearaggruppamento-descrizione-testo}
\textbf{Pre-condizioni:}
\begin{itemize}[label=$-$, noitemsep]
      \item {è in corso la definizione di una Procedura }\textit{p}
\end{itemize}
\textbf{Post-condizioni:} -

\subsection[inserisciPassoRaggruppamento]{\texorpdfstring{inserisciPassoRaggruppamento(\myuline{descrizione}: testo)}}\label{a.2-inseriscipassoraggruppamento-descrizione-testo}
\textbf{Pre-condizioni:}
\begin{itemize}[label=$-$, noitemsep]
      \item {è in corso la definizione di una Procedura }\textit{p}
\end{itemize}
\textbf{Post-condizioni:} -

\section[aggiungiIngrediente]{\texorpdfstring{aggiungiIngrediente(\myuline{ingrediente}{: Ingrediente, }\myuline{dosi}?: testo)}}\label{h.rhtfki5g8w24}
\textbf{Pre-condizioni:}
\begin{itemize}[label=$-$, noitemsep]
      \item {è in corso la definizione di una Procedura }\textit{p}
\end{itemize}
\textbf{Post-condizioni:}
\begin{itemize}[label=$-$, noitemsep]
      \item {è stata creata una istanza }\textit{i}{ di Ingrediente}
      \item \textit{i.nome }{= }\myuline{ingrediente}
      \item {{[}}\textbf{Se}{ sono specificate le }\myuline{dosi}{{]} }\textit{i.dose }{= }\myuline{dosi}
      \item \textit{i }\textbf{compone }\textit{p}
\end{itemize}

\newpage
\subsection[modificaIngrediente]{\texorpdfstring{modificaIngrediente(\myuline{ingrediente}{: Ingrediente, }\myuline{dosi}: testo)}}\label{h.btq829z3u51f}
\textbf{Pre-condizioni: }
\begin{itemize}[label=$-$, noitemsep]
      \item {è in corso la definizione di una procedura }\textit{p }{e esiste un}
            \myuline{ingrediente}{~}\textit{i}{~tale che }\textit{p}{~è }\textbf{composto}{~da }\textit{i}
\end{itemize}
\textbf{Post-condizioni: }
\begin{itemize}[label=$-$, noitemsep]
      \item \textit{i.dose }{= }\myuline{dosi}
\end{itemize}

\subsection[eliminaIngrediente]{\texorpdfstring{eliminaIngrediente(\myuline{ingrediente}: Ingrediente)}}\label{h.6moldj1jg97t}
\textbf{Pre-condizioni: }
\begin{itemize}[label=$-$, noitemsep]
      \item {è in corso la definizione di una procedura }\textit{p }{e esiste un}
            \myuline{ingrediente}{~}\textit{i}{~tale che }\textit{p}{~è }\textbf{composto}{~da }\textit{i}
\end{itemize}
\textbf{Post-condizioni: }
\begin{itemize}[label=$-$, noitemsep]
      \item {l'istanza }\myuline{ingrediente}{ viene eliminata}
\end{itemize}

\section[consultaRicettario]{\texorpdfstring{consultaRicettario()}}\label{h.w08r7f8npm2z}
\textbf{Pre-condizioni:}{ -}\newline
\textbf{Post-condizioni: }{-}\newline

\section[aggiungiPreparazioneComeIngrediente]{\texorpdfstring{aggiungiPreparazioneComeIngrediente(\myuline{preparazione}: Preparazione)}}\label{h.3umm3owhjhiq}
\textbf{Pre-condizioni:}
\begin{itemize}[label=$-$, noitemsep]
      \item {è in corso la definizione di una Procedura }\textit{p}
\end{itemize}
\textbf{Post-condizioni:}
\begin{itemize}[label=$-$, noitemsep]
      \item \myuline{preparazione}\textbf{ è ingrediente}{~in }\textit{p}
\end{itemize}

\subsection[eliminaPreparazioneComeIngrediente]{\texorpdfstring{eliminaPreparazioneComeIngrediente(\myuline{preparazione}: Preparazione)}}\label{a.1-eliminapreparazionecomeingredientepreparazione-preparazione}
\textbf{Pre-condizioni:}
\begin{itemize}[label=$-$, noitemsep]
      \item {è in corso la definizione di una Procedura }\textit{p}
\end{itemize}
\textbf{Post-condizioni:} % TODO

\section[modificaPasso]{\texorpdfstring{modificaPasso(\myuline{passo}: Passo, }\myuline{descrizione}{: testo)}}\label{h.vw6rupgvnaur}
\textbf{Pre-condizioni:}
\begin{itemize}[label=$-$, noitemsep]
      \item {è in corso la definizione di una Procedura }\textit{p}{ e }\textit{p }
            {è }\textbf{formato da }\myuline{passo}
\end{itemize}
\textbf{Post-condizioni:}
\begin{itemize}[label=$-$, noitemsep]
      \item \myuline{passo.descrizione}{~= }\myuline{descrizione}
\end{itemize}

\subsection[modificaOrdinePassi]{\texorpdfstring{modificaOrdinePassi(\myuline{ordinamento})}}\label{h.9ehcjm882e53}
\textbf{Pre-condizioni:}
\begin{itemize}[label=$-$, noitemsep]
      \item {è in corso la definizione di una Procedura }\textit{p}
\end{itemize}
\textbf{Post-condizioni: }
\begin{itemize}[label=$-$, noitemsep]
      \item {l'associazione}\textbf{~formato da}{~tra }\textit{p}{~e }\textit{i}
            {~Passi è stata modificata in accordo a }\myuline{ordinamento}
\end{itemize}

\subsection[eliminaPasso]{\texorpdfstring{eliminaPasso(\myuline{passo}: Passo)}}\label{b.1-eliminapassopasso-passo}
\textbf{Pre-condizioni:}
\begin{itemize}[label=$-$, noitemsep]
      \item {è in corso la definizione di una Procedura }\textit{p}
\end{itemize}
\textbf{Post-condizioni: } % TODO

\section[dettagliaProcedura]{\texorpdfstring{dettagliaProcedura(\myuline{annotazione}: testo, }
\myuline{porzioni}{?: numero, }\myuline{tAttvità}{?: LocalTime, } \myuline{tTotale}
{?: LocalTime, }\myuline{tUltimazione} {?: LocalTime)}}\label{h.rc4f999olfrs}
\textbf{Pre-condizioni:}
\begin{itemize}[label=$-$, noitemsep]
      \item {è in corso la definizione di una Procedura }\textit{p}
\end{itemize}
\textbf{Post-condizioni:}
\begin{itemize}[label=$-$, noitemsep]
      \item {viene creata una nuova istanza }\textit{a }{di Annotazione}
      \item \textit{a.descrizione }{= }\myuline{annotazione}
      \item \textit{p }{è }\textbf{classificata}{~da }\textit{a}
      \item {{[}}\textbf{Se}{~sono specificate le }\myuline{porzioni}{{]} }\textit{p.porzioni }{= }\myuline{porzioni}
      \item {{[}}\textbf{Se}{~è specificata una }\myuline{tAttività}{{]} }\textit{p.tempoDiAttività }{= }\myuline{tAttività}
      \item {{[}}\textbf{Se}{~è specificata una }\myuline{tTotale}{{]} }\textit{p.tempoTotale }{= }\myuline{tTotale}
      \item {{[}}\textbf{Se}{~è specificata una }\myuline{tUltimazione}{{]} }\textit{p.tempoDiUltimazione }{= }\myuline{tUltimazione}
\end{itemize}

\subsection[modificaDettagliProcedura]{\texorpdfstring{modificaDettagliProcedura(}\myuline{annotazione}{: Annotazione, }\myuline{descrizione}{: 
testo, }\myuline{porzioni?}{: numero, }\myuline{tAttvità}{?: LocalTime, }\myuline{tTotale}{?:
LocalTime, }\myuline{tUltimazione}{?: LocalTime)}}\label{h.vmiiciyams7p}
\textbf{Pre-condizioni:}
\begin{itemize}[label=$-$, noitemsep]
      \item {è in corso la definizione di una Procedura }\textit{p }{che è classificata da}{~}\myuline{annotazione}
\end{itemize}
\textbf{Post-condizioni:}
\begin{itemize}[label=$-$, noitemsep]
      \item \myuline{annotazione.tag}{ = }\myuline{descrizione}
      \item {{[}}\textbf{Se}{~sono specificate le }\myuline{porzioni}{]} \textit{p.porzioni}{ = }\myuline{porzioni}
      \item {{[}}\textbf{Se}{~è specificata una }\myuline{tAttività}{{]} }\textit{p.tempoDiAttività }{= }\myuline{tAttività}
      \item {{[}}\textbf{Se}{~è specificata una }\myuline{tTotale}{{]} }\textit{p.tempoTotale }{= }\myuline{tTotale}
      \item {{[}}\textbf{Se}{~è specificata una }\myuline{tUltimazione}{{]} }\textit{p.tempoDiUltimazione }{= }\myuline{tUltimazione}
\end{itemize}

\subsection[eliminaAnnotazione]{\texorpdfstring{eliminaAnnotazione (\myuline{annotazione}: Annotazione)}}\label{h.b1ybdk3i4rdw}
\textbf{Pre-condizioni:}
\begin{itemize}[label=$-$, noitemsep]
      \item {è in corso la definizione di una Procedura }\textit{p }{che è classificata da}{~}\myuline{annotazione}
\end{itemize}
\textbf{Post-condizioni:}
\begin{itemize}[label=$-$, noitemsep]
      \item {l'istanza }\myuline{annotazione}{ viene eliminata}
\end{itemize}

\section[pubblica]{\texorpdfstring{pubblica(\myuline{nome}?: testo)}}\label{h.b034hx47hgkh}
\textbf{Pre-condizioni:}
\begin{itemize}[label=$-$, noitemsep]
      \item {è in corso la definizione di una Procedura }\textit{p}
\end{itemize}
\textbf{Post-condizioni:} % TODO Builder pattern maybe?
\begin{itemize}[label=$-$, noitemsep]
      \item {{[}}\textbf{Se }\myuline{nome}{~}{specificato}{{]} }\textit{p.nome }{= }\myuline{nome}
      \item \textit{p.pubblicato }{= sì}
\end{itemize}

\subsection[terminaSenzaPubblicare]{\texorpdfstring{terminaSenzaPubblicare()}}\label{h.g9jb5nwngrzz}
\textbf{Pre-condizioni:}
\begin{itemize}[label=$-$, noitemsep]
      \item {è in corso la definizione di una Procedura }\textit{p}
\end{itemize}
\textbf{Post-condizioni: }{-}

\part{Progettazione}
\chapter{Domain Class Diagram}
\chapter{Design Sequence Diagram}

\end{document}